\section{Towers}
\label{sec:g-towers}
\source{CSES 1073 (Notion: Greedy)}{Medium}

\problemstatement{Cubes of sizes $k_1,\dots,k_n$ arrive in order
($n\le2\cdot10^5$). Each cube is placed on top of an existing tower whose top
cube is strictly larger, or starts a new tower. Minimise the number of
towers.}

\begin{patternbox}
``Place each item on the pile whose top is the smallest value still
larger than it'' --- patience sorting. The number of piles equals the length
of the longest non-decreasing subsequence.
\end{patternbox}

\subsection*{Approach and intuition}

Keep the multiset of tower tops. For a new cube $k$, put it on the tower with
the \emph{smallest} top $>k$; if none exists, start a new tower.

\begin{ideabox}[Lower Bound Meets Greedy]
\emph{Lower bound.} If $i<j$ and $k_i\le k_j$, cubes $i$ and $j$ can never
share a tower ($j$ would sit above $i$ but is not smaller). So a
non-decreasing subsequence of length $L$ forces at least $L$ towers.

\emph{Greedy matches it.} The greedy keeps the tops sorted. When a cube $k$
opens tower number $m$, every existing top is $\le k$; link $k$ to the current
top of tower $m-1$ (placed earlier, value $\le k$). Following these links
from the cube that opened the last tower yields a non-decreasing subsequence
whose length equals the number of towers. Hence greedy $=$ optimum.
\end{ideabox}

\subsection*{Algorithm}
\begin{algorithmic}[1]
\State $tops\gets$ empty multiset
\For{each cube $k$}
  \State $it\gets tops.\text{upper\_bound}(k)$
  \If{$it\ne tops.\text{end}()$} erase $it$ \EndIf
  \State insert $k$
\EndFor
\State \Return $|tops|$
\end{algorithmic}

\dryrun{$3,8,2,1,5$: tops $\{3\}\to\{3,8\}\to\{2,8\}$ ($2$ on $3$)
$\to\{1,8\}$ ($1$ on $2$) $\to\{1,5\}$ ($5$ on $8$). Two towers. \checkmark}

\subsection*{C++17 implementation}
\cppfile{greedy/13_towers.cpp}

\begin{complexitybox}
\textbf{Time:} $O(n\log n)$. \textbf{Space:} $O(n)$.
\end{complexitybox}

\begin{pitfallbox}
\begin{itemize}
  \item The upper cube must be \emph{strictly} smaller: equal sizes need a
        new tower, hence \texttt{upper\_bound}, not \texttt{lower\_bound}.
  \item The tops multiset is always sorted, so a plain \texttt{vector} with
        binary search also works (like LIS).
\end{itemize}
\end{pitfallbox}
